Ordinal labels encode both class identity and position. An error from class
$4$ to class $3$ differs structurally from an error from class $4$ to class
$0$, even when both are exact-class errors. This distinction is especially
important for rare endpoint classes in imbalanced ordinal settings, where
aggregate accuracy or mean error can obscure the direction of an error. A rare
upper endpoint may be assigned to an adjacent interior class, or its predicted
probability mass may be distributed across more central classes. These outcomes
have different implications for ordinal decisions, but are often grouped under
the same tail-class performance summary.

The distinction is not merely semantic. A prediction rule that preserves much
of the endpoint probability mass but makes an adjacent decision has a different
ordinal profile from one whose distribution is centered several classes inward.
For applications that use the ordered label itself, this difference affects the
size and direction of the downstream decision error. A standard class-wise
recall result cannot determine whether an endpoint failure is a local boundary
error or a broader displacement toward the interior.

We study this pattern as \emph{rare-end inward localization}. For a true upper
endpoint, inward localization is visible in discrete routing, predictive
location, and inward shrinkage of the predictive distribution. It is not
defined by low exact-class accuracy alone. Standard imbalance analyses can show
that a rare class is difficult, while leaving open whether its errors are
directionally displaced toward the interior of the ordered label space. This
distinction also separates a modest adjacent error from a larger ordinal
mislocalization. We describe inward localization in the studied imbalanced
settings; we do not assume that imbalance alone causes it.

This perspective complements, rather than replaces, standard evaluation.
Accuracy, MAE, and calibration remain useful summaries, but they do not by
themselves identify the location of a predictive distribution relative to the
true endpoint. We therefore use routing together with predictive mean and
inward shrinkage to characterize the observed phenomenon. These diagnostics
keep the ordered structure visible without treating any one scalar score as a
complete account of rare-end behavior.

Prior work has established that long-tail recognition can involve both
representation and classifier head effects. It has also developed decoupled
training, balanced classifier retraining, and logit adjustment
\cite{kang2020decoupling,alshammari2022weight,menon2021logit}; related work
also analyzes classifier and feature geometry
\cite{c2am2022angular,zhu2022balanced,yi2025geometry,wang2026alignment}.
Ordinal classification work similarly uses ordering-aware costs and
resampling to improve predictive performance under ordered labels
\cite{zhu2019smor,lazaro2023bayesian}. These directions
motivate attention to classifier heads and geometry, but they do not by
themselves characterize rare ordinal endpoints as an inward-localization
phenomenon. Our question is narrower: when a rare upper endpoint is localized
inward, how much of that localization can be altered by a classifier-head
direction intervention while the representation, class-specific scale, and bias
are held fixed?

This question does not require an assumption that every rare-end sample has the
same failure mode. A representation may already place some endpoint samples
closer to an interior class, whereas other samples may retain endpoint-like
geometry before the original head maps them inward. The distinction motivates a
decomposition and a constrained intervention, while leaving the centroid-based
representation diagnostic descriptive rather than causal.

We address this question through a mechanism analysis rather than a new
classifier proposal. We first diagnose rare-end inward localization and use
historical centroid-based analyses to distinguish representation-inward samples
from samples that retain rare-end-like geometry but are mapped inward by the
original head. We then compare the original frozen head with a direction-only
head adaptation that preserves the original row norms and biases. The primary
test uses independently trained CE- and RPS-based backbone seeds in
RetinaMNIST and Solar. Across both datasets and objectives, all four
confirmatory backbone seeds per setting reduce rare-end localization error under
the direction-only intervention. Endpoint-specific representation geometry adds
explanatory value strongly in Solar, is partial in RetinaMNIST CE, and is not
supported by the preregistered criterion in RetinaMNIST RPS. We also use
controlled RetinaMNIST studies to examine the adaptation signal and the
boundary of a rare-support interpretation.

The cross-setting replication and the geometry analysis answer different
questions. The former asks whether the direction-only response recurs across
independently trained backbones within each tested setting. The latter asks
whether one endpoint-specific representation measure adds explanatory value
after accounting for original-head state and generic geometry. Reporting both
the strong and mixed outcomes prevents the more stable intervention response
from being mistaken for a universal account of recoverability.

Our contributions are:
\begin{itemize}
    \item We define and characterize rare-end inward localization using
    ordinal routing, predictive location, and shrinkage in the studied
    imbalanced ordinal settings.
    \item We show that a fixed-norm, fixed-bias direction-only classifier-head
    intervention reproducibly improves rare-end localization across
    RetinaMNIST and Solar under CE- and RPS-trained representations, with the
    effect sign replicated in all four confirmatory backbone seeds per setting.
    \item We bound the mechanism interpretation: the historical
    representation/head decomposition is descriptive, endpoint-specific
    geometry is setting-dependent, and the controlled factor and support
    studies do not justify a universal rebalancing or imbalance-causality
    claim.
\end{itemize}
